\Subhead{1.8.2}{Спектральные методы оценки устойчивости}
\paragraph{Периодический режим: Флоке и Пуанкаре.}
Вдоль решения с периодом $T$ вариационное уравнение имеет вид
\begin{equation}
 \dot{\bm\xi}=A(t)\bm\xi,\quad A(t+T)=A(t),\quad
 \dot\Phi=A(t)\Phi,\quad\Phi(0)=I.
\end{equation}
Собственные значения $\Phi(T)$ --- мультипликаторы Флоке $\rho_j$. После $n$ периодов возмущение меняется как $\rho_j^n$; показатели удовлетворяют $\operatorname{Re}\lambda_j=T^{-1}\ln|\rho_j|$. Для вынужденно-периодического режима асимптотическая устойчивость требует $|\rho_j|<1$. У автономного цикла одна мультипликаторная величина равна $1$ из-за сдвига времени; устойчивость орбиты требует, чтобы остальные лежали внутри единичного круга. На поперечном сечении Пуанкаре это нейтральное направление отсутствует.

\paragraph{Частотная проверка обратной связи.}
Для скалярной отрицательной обратной связи с разомкнутой передаточной функцией $L(s)$ характеристическое уравнение замкнутого контура $1+L(s)=0$. Временной метод проверяет затухание отклика, спектральный --- положение полюсов $s$. Критерий Найквиста позволяет определить число неустойчивых корней по частотному годографу $L(\ii\omega)$ и известному числу правых полюсов разомкнутого контура.

Зададим соглашение однозначно: контур в $s$-плоскости обходит правую полуплоскость \emph{по часовой стрелке}; $N_{\mathrm{cw}}$ --- результирующее число обходов точки $-1$ его образом $L(s)$ по часовой стрелке. Принцип аргумента даёт
\begin{equation}
 Z=P+N_{\mathrm{cw}},
\end{equation}
где $P$ --- число правых полюсов $L$, $Z$ --- число правых нулей $1+L$. Для устойчивости нужно $Z=0$: при $P=0$ годограф не должен охватывать $-1$, а при $P>0$ нужны $P$ обходов против часовой стрелки. На мнимой оси контур деформируют; проход через критическую точку соответствует границе устойчивости. Должны отсутствовать скрытые неустойчивые сокращения полюсов и нулей. Для положительной связи $1-L=0$ критическая точка меняется на $+1$.

АЧХ без ФЧХ недостаточно: при той же амплитуде сдвиг фазы может изменить знак обратной связи. Условие единичного усиления на одной частоте определяет возможную границу, но не доказывает устойчивость всего контура. Спектр \emph{самих наблюдаемых колебаний} тоже не заменяет спектр вариационного оператора: узкая линия не доказывает устойчивости режима.
